We hope to introduce the reader to the main ideas and methods that will be used throughout this work.

\begin{definition}[Chain rule for densities]
    For random variables $Y_1,\dots,Y_n$,
    \[
    f_{Y_1,\dots,Y_n}(y_1,\dots,y_n)
    = f_{Y_1}(y_1)
    \prod_{j=2}^n f_{Y_j \mid Y_1,\dots,Y_{j-1}}(y_j \mid y_1,\dots,y_{j-1}).
    \]
    This identity, called the \emph{prediction decomposition}, expresses the joint density as the product of one-step-ahead conditional densities.
\end{definition}

\subsection{AR(1) Gaussian Process}
The following process is an autoregressive model of order~1 (AR(1)):
$$Y_j = \mu + \alpha(Y_{j-1}-\mu) + \varepsilon_j, \quad \varepsilon_j \sim \mathcal{N}(0,\sigma^2),$$

whuch satisfies the following conditional distribution:
$$Y_j \mid Y_{j-1}=y_{j-1} \sim \mathcal{N}\{\mu+\alpha(y_{j-1}-\mu),\,\sigma^2\}.$$

\subsection{Likelihood}
As a preliminary note, we define the likelihood function $L(\textbf{y} | \theta)$ as $\mathbb{P}(\text{seeing data} | \theta) = \mathbb{P}(y_1, \ldots, y_n | \theta)$ which
`splits' due to the property of independence. More formally,

\begin{definition}[Likelihood]
    Let $\textbf{y} = (y_1, \ldots, y_n)$ be i.i.d. samples from a probability mass function (pmf) $p_{Y}(t | \theta)$ (if $Y$ is discrete) or probability density
    function (pdf) $f_{Y}(t | \theta)$ (if $Y$ is continuous) parameterized (with parameter) by $\theta \in \Theta$. We define
    the likelihood of $\textbf{y}$ given $\theta$ to be the “probability” of observing \textbf{y} if the true parameter is $\theta$.
    If $Y$ is discrete,
    $$L(\textbf{y}| \theta) = \prod_{i=1}^{n}p_{Y}(y_i | \theta)$$

    and if $Y$ is continuous,
    $$L(\textbf{y}| \theta) = \prod_{i=1}^{n}f_{Y}(y_i | \theta).$$

\end{definition}

Hence, the likelihood function for data $y_0,\ldots,y_n$ from an AR(1) Gaussian process with known initial value $y_0$ is given by
\[
L(\mu,\alpha,\sigma^2) = \prod_{j=1}^n \frac{1}{\sqrt{2\pi\sigma^2}} \exp\left\{-\frac{1}{2\sigma^2}[y_j - \mu - \alpha(y_{j-1}-\mu)]^2\right\}.
\]

\begin{remark}
    The log-likelihood function is often more convenient to work with, as it transforms products into sums. The log-likelihood for the AR(1) Gaussian process is
    \[
    \ell(\mu,\alpha,\sigma^2) = -\frac{n}{2}\log(2\pi\sigma^2) - \frac{1}{2\sigma^2}\sum_{j=1}^n [y_j - \mu - \alpha(y_{j-1}-\mu)]^2.
    \]
\end{remark}